\chapter{Diff de hunts}
\label{ch:hunt-diff}

\code{diff\_hunts} responde a la pregunta ``¿qué cambió entre dos
puntos en el tiempo en el conjunto de caminos que satisfacen
esta hipótesis?'' Es la base del panel ``qué hay nuevo desde mi
último chequeo'' en la UI y de la automatización programada que
alerta sobre cadenas de ataque recién aparecidas.

\section{Variante clásica por tiempo de evento}

El algoritmo base corre la hipótesis dos veces contra el grafo,
acotando la ventana temporal a distintos upper-bounds, y diffea
los conjuntos de caminos.

\begin{algorithm}
\caption{\textsc{Diff-Hunts}}\label{alg:diff-hunts}
\KwIn{grafo $G$, hipótesis $H$, timestamp baseline $t_{\text{base}}$,
      timestamp actual $t_{\text{cur}}$ con $t_{\text{base}} < t_{\text{cur}}$}
\KwOut{(agregados: $\mathrm{List}[\pi]$, quitados: $\mathrm{List}[\pi]$)}
$\Pi_{\text{base}} \gets $ \textsc{Search-Temporal-Pattern}($G, H, (0, t_{\text{base}}), \infty$)\;
$\Pi_{\text{cur}} \gets $ \textsc{Search-Temporal-Pattern}($G, H, (0, t_{\text{cur}}), \infty$)\;
$U_{\text{base}} \gets $ \textsc{Score-And-Paginate}($\Pi_{\text{base}}$, 0, $\infty$, None, ByPath).únicos\;
$U_{\text{cur}} \gets $ idem para $\Pi_{\text{cur}}$\;
$K_{\text{base}} \gets $ conjunto de claves de camino en $U_{\text{base}}$\;
$K_{\text{cur}} \gets $ idem para $U_{\text{cur}}$\;
agregados $\gets U_{\text{cur}}$ filtrados por clave $\notin K_{\text{base}}$\;
quitados $\gets U_{\text{base}}$ filtrados por clave $\notin K_{\text{cur}}$\;
\Return (agregados, quitados)\;
\end{algorithm}

\begin{complexity}
Dos ejecuciones de hunt más una pasada de diferencia de
conjuntos. Si $T_{\text{match}}$ es el costo del matcher y $P$ es
la cantidad de caminos devueltos por cualquiera de los hunts, el
total es $2 T_{\text{match}} + \Ord{P}$. Espacio $\Ord{P}$ por
encima del scratch propio del matcher.
\end{complexity}

\section{Trampa semántica: datos que llegan tarde}

La variante clásica diffea por \emph{tiempo de evento}, no por
\emph{tiempo de conocimiento}: una arista cuyo timestamp de
evento es del pasado pero que fue ingerida después de
$t_{\text{base}}$ aparecerá tanto en $\Pi_{\text{base}}$ como en
$\Pi_{\text{cur}}$, porque el upper-bound de la ventana temporal
no filtra basado en cuándo el motor se enteró de la arista. Esta
es la motivación original del trabajo MVCC P3
(Capítulo~\ref{ch:mvcc}): un analista que cerró un ticket en
$t_{\text{base}}$ igual puede querer saber que una ingesta
post-mortem reveló aristas nuevas dentro de la ventana que ya
consideraba ``cerrada''.

\section{La variante por snapshot (diferida)}

Una vez que aterrice la extensión de matcher P3
(\textsc{Search-Temporal-Pattern-At}), \code{diff\_by\_snapshot}
se vuelve la alternativa más rica:

\begin{algorithm}
\caption{\textsc{Diff-By-Snapshot} (diferida)}\label{alg:diff-snapshot}
\KwIn{grafo $G$, hipótesis $H$, ventana de evento $(t_0, t_1)$,
      snapshot baseline $\iota_{\text{base}}$, snapshot actual
      $\iota_{\text{cur}}$}
\KwOut{listas de caminos (agregados, quitados)}
$\Pi_{\text{base}} \gets $ \textsc{Search-Temporal-Pattern-At}($G, H, (t_0, t_1), \iota_{\text{base}}$)\;
$\Pi_{\text{cur}} \gets $ \textsc{Search-Temporal-Pattern-At}($G, H, (t_0, t_1), \iota_{\text{cur}}$)\;
\ldots set-diff como en \textsc{Diff-Hunts} \ldots
\end{algorithm}

Distinción clave: la ventana de evento $(t_0, t_1)$ es fija, y lo
que varía es el snapshot de conocimiento $\iota$. El motor
responde ``¿qué caminos nuevos en esta ventana de evento me
enteré entre $\iota_{\text{base}}$ y $\iota_{\text{cur}}$?'' ---
la semántica apropiada de auditoría para enriquecimiento que
llega tarde, detonaciones VirusTotal que vuelven después de
ingerido el evento, overlays de threat intel que llegaron más
tarde, etc.

\begin{complexity}
Igual que la variante clásica: $2 T_{\text{match}} + \Ord{P}$.
El delta de \emph{valor} es semántico, no de performance.
\end{complexity}

\section{Validación de entrada}

\begin{invariant}[Orden de la ventana]
\code{diff\_hunts} rechaza $t_{\text{base}} > t_{\text{cur}}$ con
\code{ApiError::InvalidInput}. Una ventana invertida produciría
un par (agregados, quitados) cuya semántica es responsabilidad
del usuario de interpretar; el motor prefiere fallar ruidosamente.
\end{invariant}

Este invariante lo ejercita el arnés de paridad
(\code{diff\_hunts\_rejects\_inverted\_window} en
\filepath{graph\_hunter\_api/tests/parity/scenarios.rs}).

\begin{inthecode}
\filepath{graph\_hunter\_core/src/analytics.rs:1342} ---
\code{diff\_hunts} (clásico).\\
\filepath{graph\_hunter\_core/src/relation.rs:83} --- el campo
\code{ingested\_at\_delta} que \code{diff\_by\_snapshot}
eventualmente consultará.
\end{inthecode}
